\begin{textsample}{POS Dim 1 – vem\_ed\_28 – Score 126.00 – vem\_ed\_28/t149.txt}  \label{ex:f1_pos_001}
Escrituras narrativas: \textbf{diferenças} em \textbf{diálogo} “Torna-se \textbf{cada} vez mais raro o \textbf{encontro} com \textbf{pessoas} \textbf{que} \textbf{sabem} narrar alguma \textbf{coisa} direito” (Benjamin, 1980, p. 57).
Tomo essa citação do texto originalmente de 1936 de Walter Benjamin como provocação para uma urgência e \textbf{necessidade} de se recuperar uma \textbf{capacidade} \textbf{essencial} do \textbf{ser} \textbf{humano}: a de narrar experiências.
Para Benjamin, a narração \textbf{era} uma prática social fundamental para a transmissão de sabedoria e \textbf{vivências}, além de atuar como resistência à alienação provocada pela lógica da \textbf{produção} em massa.
Mais recentemente, o filósofo sul-coreano Byung-Chul Han (2023) alertou para a \textbf{crise} da narração, \textbf{observando} \textbf{que}, na sociedade \textbf{contemporânea}, marcada pela competição e pela superficialidade das \textbf{interações}, a \textbf{verdadeira} narração – aquela \textbf{que} \textbf{gera} \textbf{empatia} e \textbf{compreensão} – tornou-se rara.
Embora em um \textbf{nível} mais básico e introdutório, condizente com o ensino de espanhol para brasileiros, as \textbf{reflexões} de Benjamin e Han inspiram atividades, \textbf{conteúdos} e abordagens do Projeto Colaborativo Internacional (PCI / Cesu) “Escrituras Narrativas”, desenvolvido com a Uniminuto (Colômbia).
No primeiro semestre de 2025, o projeto alcança sua quarta edição.
Nos cursos de espanhol para fins \textbf{específicos} da Fatec Ipiranga, enfatizo a \textbf{importância} de \textbf{saber} \textbf{contar}, uma \textbf{habilidade} mais \textbf{ampla} \textbf{que} dar respostas curtas e informações pontuais.
Narrar bem envolve diversos aspectos da vida cotidiana, tanto no âmbito pessoal \textbf{quanto} no \textbf{profissional}. \textbf{continuação} Dessa \textbf{forma}, a \textbf{estrutura} do ensino de espanhol para o curso de Gestão Comercial se organiza em três \textbf{momentos}: inicialmente, os alunos aprendem a \textbf{contar} quem \textbf{são}; em seguida, a descrever o \textbf{que} fazem; e, por fim, a \textbf{compartilhar} suas experiências passadas e planos \textbf{futuros}.
O PCI / Cesu insere-se precisamente nesse \textbf{terceiro} \textbf{momento}.
O projeto envolve alunos colombianos do curso \textbf{superior} de Tecnologia em Realização Audiovisual e estudantes da Fatec Ipiranga.
Com exceção da segunda edição, realizada no segundo semestre de 2023 com a \textbf{turma} de Gestão de Eventos, \textbf{todas} as demais edições \textbf{foram} \textbf{implementadas} com \textbf{turmas} do \textbf{terceiro} semestre de Gestão Comercial.
A metodologia adotada baseia-se em \textbf{princípios} das metodologias \textbf{ativas}, da aprendizagem \textbf{significativa} (Carril, Natário e Zoccal, 2017) e da \textbf{pedagogia} da autonomia (Freire, 1996), \textbf{utilizando} a aprendizagem \textbf{baseada} em projetos (Bender, 2014) como meio de \textbf{promover} a \textbf{integração} entre estudantes de diferentes instituições, cursos, \textbf{idiomas}, países e \textbf{contextos} \textbf{profissionais}.
Para \textbf{alinhar} tantos \textbf{elementos} distintos e encontrar um \textbf{ponto} de convergência benéfico para ambos os grupos, tornou-se fundamental o trabalho intenso e \textbf{contínuo} com a professora Magda Lorena Parodi Ruiz, parceira internacional com quem desenvolvo \textbf{todas} as edições do projeto.
Quando um projeto passa por várias edições com os mesmos parceiros, a experiência \textbf{acumulada} se \textbf{torna} um \textbf{fator} crucial para sua melhoria \textbf{contínua}.
O \textbf{momento} de avaliação pós-projeto \textbf{é} \textbf{essencial} para \textbf{identificar} \textbf{pontos} a \textbf{serem} ajustados e resgatar práticas \textbf{bem-sucedidas}.
Nesse \textbf{contexto}, um dos aspectos \textbf{analisados} \textbf{foi} a duração do projeto.
Experimentamos diferentes formatos, \textbf{variando} entre oito, seis e cinco semanas.
A edição mais curta \textbf{resultou} em \textbf{processos} acelerados e \textbf{menor} tempo para \textbf{reflexão} e elaboração.
Embora o resultado não tenha \textbf{sido} comprometido, \textbf{foi} necessário um esforço extra, realizado fora dos \textbf{momentos} planejados, o \textbf{que} poderia \textbf{ser} evitado com um \textbf{cronograma} mais equilibrado.
Assim, embora o projeto possa ocorrer em um período \textbf{reduzido} (cinco semanas), a configuração \textbf{ideal} para o PCI / Cesu aqui relatado envolve \textbf{encontros} \textbf{síncronos} e atividades assíncronas distribuídas ao \textbf{longo} de oito semanas.
No \textbf{entanto}, essa duração mais \textbf{longa} também apresenta desafios, como a conciliação dos \textbf{programas} acadêmicos das instituições envolvidas e a sobrecarga dos alunos com outras atividades \textbf{curriculares}. \textbf{continuação} Portanto, uma avaliação prévia do estado das \textbf{turmas} a \textbf{cada} semestre \textbf{é} \textbf{essencial} para ajustar a carga \textbf{horária} e os objetivos do projeto de \textbf{forma} realista.
Em três edições, o trabalho \textbf{foi} realizado com alunos do \textbf{terceiro} semestre de Gestão Comercial e, na segunda edição, com alunos do \textbf{quinto} semestre de Gestão de Eventos.
O objetivo \textbf{comum} às duas \textbf{disciplinas} \textbf{foi} \textbf{consolidar} e aplicar \textbf{conhecimentos} do ensino de espanhol para fins \textbf{específicos}, \textbf{expandindo} sua utilização para diferentes \textbf{contextos} e depois retornando a esse \textbf{foco} com novas \textbf{perspectivas}.
Como o ensino de espanhol para fins \textbf{específicos} para cursos de gestão \textbf{permite} \textbf{alinhar} objetivos acadêmicos e \textbf{profissionais}, as configurações do projeto \textbf{foram} semelhantes nas duas primeiras edições, com \textbf{mudanças} mais \textbf{significativas} nas duas seguintes devido à participação de uma \textbf{disciplina} diferente por \textbf{parte} dos parceiros colombianos.
Nas primeiras edições, a professora Parodi Ruiz trabalhou com seus alunos na \textbf{disciplina} “Escrituras Narrativas”, cujo objetivo \textbf{era} desenvolver roteiros cinematográficos consistentes e esteticamente elaborados, \textbf{seguindo} a \textbf{estrutura} clássica de roteiro cinematográfico (\textbf{início}, \textbf{conflito} e clímax).
Durante o projeto, os estudantes colombianos \textbf{compartilhavam} imagens de suas \textbf{produções} com os colegas brasileiros, \textbf{que} \textbf{analisavam} e propunham narrativas para essas imagens.
Esse exercício \textbf{gerou} debates aprofundados sobre o \textbf{papel} da imagem na percepção e \textbf{compreensão} da \textbf{realidade}, além de \textbf{conexões} com o \textbf{storytelling} aplicado à \textbf{construção} de marcas, \textbf{produtos} e serviços.
Na \textbf{terceira} e quarta edições, o projeto passou a envolver alunos colombianos do primeiro semestre na \textbf{disciplina} “Taller Creativo” (Oficina de criatividade).
Diferentemente das edições \textbf{anteriores}, em \textbf{que} os alunos trabalhavam na \textbf{construção} de narrativas audiovisuais completas, essa \textbf{disciplina} tem como \textbf{foco} a exploração da criatividade por meio da experimentação visual com cores, texturas e \textbf{emoções}.
Consequentemente, o título do projeto \textbf{foi} \textbf{ampliado} para “Escrituras Narrativas a través de la imagen”. \textbf{continuação} Os estudantes brasileiros passaram a interpretar essas sequências imagéticas e \textbf{criar} textos narrativos \textbf{que}, ao contrário das edições \textbf{anteriores}, não \textbf{eram} \textbf{comparados} com um roteiro final de curta-metragem.
Ainda assim, essa nova abordagem proporcionou discussões e aprendizagens igualmente valiosas.
O \textbf{idioma} predominante em \textbf{todas} as atividades \textbf{foi} o espanhol, com o português \textbf{sendo} \textbf{utilizado} apenas em \textbf{momentos} de dificuldade na comunicação.
Como \textbf{estratégia} para garantir a \textbf{compreensão} \textbf{mútua}, \textbf{foi} \textbf{estabelecida} a regra de “hablar despacio y articulando bien las palabras” durante os \textbf{encontros} \textbf{síncronos}.
Quando necessário, os professores atuavam como mediadores ou intérpretes.
Além disso, o chat escrito \textbf{foi} \textbf{utilizado} para esclarecimento de vocabulário e expressões mais complexas.
A organização do projeto \textbf{seguiu} os \textbf{princípios} da aprendizagem \textbf{baseada} em projetos, conforme descritos por Bender (2014): \textbf{foco} em problemas reais e desafiadores, aprendizagem \textbf{ativa} e prática, trabalho \textbf{colaborativo}, autonomia do aluno, \textbf{integração} de diferentes \textbf{disciplinas}, \textbf{desenvolvimento} de \textbf{competências}, avaliação \textbf{contínua} e \textbf{produção} de um \textbf{produto} \textbf{significativo}.
Nas duas primeiras edições, a entrega incluiu vídeos combinando sequências de imagens \textbf{criadas} pelos alunos colombianos com narrativas textuais e em áudio \textbf{desenvolvidas} pelos alunos brasileiros.
Esses \textbf{materiais} \textbf{foram} \textbf{utilizados} tanto como objetos de aprendizagem no ensino de espanhol \textbf{quanto} como \textbf{apoio} na pesquisa para a \textbf{produção} de curtas-metragens dos estudantes colombianos.
Ainda \textbf{que} as edições realizadas com alunos colombianos do primeiro semestre não tenham resultado em curtas-metragens, uma \textbf{contribuição} de destaque diferente deu-se entre os alunos ingressantes da Uniminuto e os alunos veteranos da Fatec: o compartilhamento de dúvidas e experiências da \textbf{rotina} acadêmica.
Dessa \textbf{forma}, a \textbf{proposta} de internacionalização para cursos da Fatec \textbf{proposta} pelos PCIs, cumpre satisfatoriamente sua missão e vai além, alcançando além dos já experimentados benefícios no âmbito das experiências \textbf{linguísticas} e da \textbf{interculturalidade}, uma intensa e \textbf{relevante} \textbf{troca} de \textbf{saberes} e experiências. \textbf{Referências} BENDER, W.
N.
Aprendizagem \textbf{baseada} em projetos: educação diferenciada para o \textbf{século} XXI.
Porto Alegre: Penso, 2014.
BENJAMIN, W.
O narrador - Observações sobre a obra de Nikolai Leskow.
In: BENJAMIN, W.; HORKHEIMER, M.; ADORNO, T.; HABERMAS, J.
Textos \textbf{escolhidos}.
São Paulo: Abril Cultural, 1980. (p. 57-74) CARRIL, M. da G.
P.; NATÁRIO, E.
G.; ZOCCAL, S.
Considerações sobre aprendizagem \textbf{significativa} a \textbf{partir} da visão de Freire e Ausubel: uma \textbf{reflexão} teórica. e-Mosaicos, v. 6, n. 13, dez. 2017.
FREIRE, P.
Pedagogia da autonomia: \textbf{saberes} necessários à prática \textbf{educativa}. 25 ed.
São Paulo: Paz e Terra, 1996.
HAN, B.-C.
A \textbf{crise} da narração.
Rio de Janeiro: Vozes, 2023.

% matched lemmas: acumular, alinhar, ampliar, amplo, analisar, anterior, apoio, ativo, basear, bem-sucedido, cada, capacidade, coisa, colaborativo, comparar, compartilhar, competência, compreensão, comum, conexão, conflito, conhecimento, consolidar, construção, contar, contemporâneo, contexto, conteúdo, continuação, contribuição, contínuo, criar, crise, cronograma, curricular, desenvolvimento, desenvolvir, diferença, disciplina, diálogo, educativo, elemento, emoção, empatia, encontro, entanto, escolher, específico, essencial, estabelecer, estratégia, estrutura, expandir, fator, foco, forma, futuro, gerar, habilidade, horário, humano, ideal, identificar, idioma, implementar, importância, integração, interação, interculturalidade, início, linguístico, longo, material, momento, mudança, mútuo, necessidade, nível, observar, papel, parte, partir, pedagogia, pequeno, permitir, perspectiva, pessoa, ponto, princípio, processo, produto, produção, profissional, programa, promover, propor, proposta, quanto, que, quinto, realidade, reduzir, referência, reflexão, relevante, resultar, rotina, saber, seguir, ser, significativo, storytelling, superior, século, síncrono, terceiro, todo, tornar, troca, turma, utilizar, variar, verdadeiro, vivência
\end{textsample}
